\documentclass[11pt]{article}

\def\chapitre{2}
\def\pagetitle{Polynômes.}

\input{setup.tex}

\begin{document}

\input{title.tex}

\thispagestyle{fancy}

\section{Quelques révisions.}

\begin{ex}{Un polynôme périodique est constant.}{}
    \begin{enumerate}
        \item Montrer qu'une fonction continue sur $\R$ et périodique est bornée.
        \item Montrer de deux manières qu'une fonction polynomiale périodique est constante.
    \end{enumerate}
    \tcblower
    \boxed{1.} Soit $f$ continue sur $\R$ et périodique de période $T$. Alors $f$ est bornée sur $[0,T]$ par TBA car continue.\\
    Donc $\exists M \in \R_+, ~ \forall x \in [0,T], ~ |f(x)| \leq M$, or $\forall x \in \R, ~ \exists n \in \Z, ~ x-nT\in[0,T]$.\\
    Donc par $T$-périodicité de $f$, $|f(x)|=|f(x-nT)|\leq M$.\\
    \boxed{2.} Soit $P\in\R[X]$ périodique, on a donc $P$ continu et périodique donc borné d'après $1.$\\
    Or les seuls polynômes bornés sont les polynômes constants (par limite en $\infty$) donc $P$ est constant.
\end{ex}

\begin{ex}{Une caractérisation des fonctions polynomiales.}{}
    \begin{enumerate}
        \item Montrer que $f$ est un polynôme de degré inférieur à $k$ ssi $f$ est dérivable $(k+1)$ fois et $f^{(k+1)}=0$.
        \item Donner une caractérisation des fonctions polynomiales.
        \item Montrer que $e^x$, $\sin x$ ne sont pas des fonctions polynomiales.
    \end{enumerate}
    \tcblower
    \boxed{1.\ra} Soit $f$ un polynôme de degré inférieur à $k$, on a immédiatement que $f^{(k+1)}=0$.\\
    \boxed{1.\la} Soit $f$ dérivable tel que $f^{(k+1)}=0$, alors en primitivant $k+1$ fois, $\deg f \leq k - 1$.\\
    \boxed{2.} Une fonction est polynômiale ssi l'une de ses dérivées est nulle.\\
    \boxed{3.} Elles ne respectent pas la caractérisation.
\end{ex}

\begin{ex}{Des équations différentielles.}{}
    \begin{enumerate}
        \item Résoudre les équations différentielles $y'=y$ et $xy'=y$.
        \item Quelles sont les solutions polynomiales parmis leurs solutions ?
        \item Résoudre dans $\K[X]$ les équations $P'=P$, puis $XP'=P$ sans utiliser les questions précédentes.
    \end{enumerate}
    \tcblower
    \boxed{1.} On a $y'=y$ ssi $y\in\{x\mapsto \l e^x, ~ \l\in\R\}$; on a $xy'=y$ ssi $y'-\frac{1}{x}y=0$ ssi $y\in \{x\mapsto \l x, ~ \l \in \R\}$.\\
    \boxed{2.} Pour la première, seulement $0$, pour la seconde, toutes.\\
    \boxed{3.} On a $P'=P$ si $\deg(P)-1=\deg(P)$, seul $P=0$ convient.\\
    On a $XP'=P$ ssi $\sum_{k=0}^d ka_k X^k=\sum_{k=0}^d a_k X^k$ ssi $\forall k \neq 1, ~ ka_k=a_k=0$, et $\exists \l \in \R, ~ a_1=\l$. 
\end{ex}

\vspace*{-0.8cm}

\begin{ex}{Anneau intègre.}{}
    \begin{enumerate}
        \item Vérifier que $\K[X]$ est intègre.
        \item Vérifier que $(\F(\R,\R), +, \times)$ est un anneau non-intègre.
        \item Soit $(A,+,\times)$ un anneau commutatif. À quelle condition l'anneau $(A[X],+,\times)$ est-il intègre ?
    \end{enumerate}
    \tcblower
    \boxed{1.} Soient $P,Q\in\K[X]$ non nuls. Alors $\deg(PQ)=\deg(P)+\deg{Q}\in\N$ donc $PQ\neq0$.\\
    \boxed{2.} Soit $f$ nulle partout sauf en $0$, et $g$ nulle partout sauf en $1$. Leur produit est $\zero$, mais $f,g\neq\zero$.\\ 
    \boxed{3.} Soient $P,Q\in\K[X]$ non nuls. Alors $PQ=a_0b_0+...+a_Nb_{N'}X^{N+N'}$.\\
    Si $(A,+,\times)$ est intègre, alors $a_N\neq0$ et $b_{N'}\neq0$ implique $a_Nb_{N'}\neq0$.\\
    Donc $P\neq0$ et $Q\neq0$ implique $PQ\neq0$. $(A,+,\times)$ intègre $\iff$ $(A[X],+,\times)$ intègre.
\end{ex}

\vspace*{-0.8cm}

\begin{ex}{}{}
    Montrer que les polynômes de $\R[X]$ de degré impair ont forcément une racine réelle.
    \tcblower
    Soit $P\in\R[X]$, qu'on suppose de degré impair $n\in\N$, on note $a_n\neq0$ son coefficient dominant.\\
    Alors $P(x)\sim_{\infty} a_nx^n$, nous allons supposer $a_n>0$, le cas négatif est similaire.
    \begin{equation*}
        \lim_{x\to-\infty} P(x) = \lim_{x\to-\infty} a_nx^n = -\infty; \et \lim_{x\to+\infty}P(x) = \lim_{x\to+\infty} a_nx^n = +\infty.
    \end{equation*}
    Par définition de la limite, $\exists a < 0, ~ \forall x < a, ~ P(x) < 0$ et $\exists b > 0, ~ \forall x > b, ~ P(x) > 0$.\\
    Puisque $P$ est un polynôme, il est continu sur $[a,b]$, on lui applique donc le TVI : $\exists x_0 \in [a,b], ~ P(x)=0$. 
\end{ex}

\vspace*{-0.8cm}

\begin{ex}{Polynômes scindés.}{}
    \begin{enumerate}
        \item Montrer que si $P\in\R[X]$ est scindé à racines simples, alors $P'$ aussi.
        \item Montrer que si $P\in\R[X]$ est scindé, alors $P'$ aussi.
    \end{enumerate}
    \tcblower
    \boxed{1.} Soit $P\in\R[X]$ scindé à racines simples : $P=\l(X-\a_1)...(X-\a_p)$. On suppose que les $\a_i$ sont croissants.\\
    Par Rolle, $P'$ s'annule une fois dans tous les $]a_i,a_{i+1}[$ car $P$ est continue et dérivable sur $[a_i,a_{i+1}]$.\\
    Donc $P'$ s'annule $p-1$ fois, et est de degré $p-1$, il est scindé à racines simples.\\
    \boxed{2.} Soit $P\in\R[X]$ scindé de degré $n=n_1+...+n_p$ : $P=\l(X-\a_1)^{n_1}...(X-\a_p)^{n_p}$. On suppose les $\a_i$ croissants.\\
    En appliquant (1), on obtient $(b_1,...b_{p-1})$, racines distinctes de $P'$ (et distinctes des $(\a_i)$).\\
    Rappelons qu'une racine de $P$ de multiplicité $n$ est racine de $P'$ de multiplicité $n-1$, alors :\vspace*{-0.2cm}
    \begin{equation*}
        \exists Q \in \R[X], ~ P'(X) = (X-b_1)...(X-b_{p-1})(X-\a_1)^{n_1-1}...(X-\a_p)^{n_p-1}Q(X).
    \end{equation*}
    Or $P'$ est de degré $n-1$, donc $n-1=n-1+\deg(Q)$, donc $Q$ est constant, donc $P'$ est scindé.
\end{ex}

\begin{ex}{Relations coefficients racines.}{}
    Soient $x,y,z$ trois complexes non nuls tels que $x+y+z=0$ et $\frac{1}{x}+\frac{1}{y}+\frac{1}{z}=0$.\\
    Montrer que $|x|=|y|=|z|$.
    \tcblower
    On pose $P=(X-x)(X-y)(X-z)$, alors $P=X^3-(x+y+z)X^2+(xy+xz+yz)X-xyz$. D'après les hypothèses :
    \begin{equation*}
        P=X^3-xyz ~ \nt{car} ~ \frac{1}{x}+\frac{1}{y}=\frac{1}{z}=\frac{yz+xz+xy}{xyz} = 0 \ra yz+xz+xy=0.
    \end{equation*}
    Donc $P(X)=X^3-xyz$, avec $x,y,z$ racines de $P$, donc $x^3=y^3=z^3=xyz$.\\
    On en conclut que $|x|^3=|y|^3=|z|^3$ puis que $|x|=|y|=|z|$.
\end{ex}

\section{Arithmétique dans \texorpdfstring{$\K[X]$}{Lg}.}

\subsection{PGCD, PPCM de deux polynômes.}

\begin{prop}{Division euclidienne.}{}
    \begin{equation*}
        \forall (A,B) \in \K[X]^2, ~ B\neq0 \ra \exists!(Q,R) \in \K[X]^2 ~ A = BQ + R ~\et~ \deg(R) < \deg(B).
    \end{equation*}
    \tcblower
    \bf{Unicité.} Supposons que $(Q,R),(Q',R')\in\K[X]$ conviennent.\\
    Alors $B(Q-Q')=R'-R$ or $\deg(R'-R)\leq\max(\deg(R),\deg(R'))\leq\deg(B)$.\\
    De plus, si $Q-Q'\neq0$, alors $\deg(B(Q-Q'))=\deg(B)+\deg(Q-Q')\geq\deg(B)$, absurde.\\
    Donc $Q-Q'=0$, puis $R'-R=0$.\\
    \bf{Existence.} Si $\deg(A)<\deg(B)$, on a le couple $(0,A)$ qui convient. Sinon par récurrence :\\
    \bf{Initialisation.} Pour $A=a\in\K$, $B=b\in\K$, on a $A=\frac{a}{b}\cdot B$.\\
    \bf{Hérédité.} Soit $n\in\N$ tel que pour $k\leq n$, la propriété soit vraie. Soit $A=a_0+...+a_{n+1}X^{n+1}$.\\
    Soit $B=b_0+...+b_pX^p$ où $p=\deg(B)\leq n+1$. On pose $A'=A-\frac{a_{n+1}}{b_p}X^{n+1-p}B\in\K[X]$.\\
    Évident : $\deg(A')\leq n+1$, et $a'_{n+1}=a_{n+1}-\frac{a_{n+1}}{b_p}b_p=0$ donc $\deg(A')\leq n$.\\
    Par hyp. : $\exists Q',R'\in\K[X], ~ A'=BQ'+R'$ avec $\deg(R')<\deg(B)$ d'où $A=B(Q'+\frac{a_{n+1}}{b_p}X^{n+1-p})+R'$.
\end{prop}

\vspace*{-0.7cm}

\begin{exi}{}{}
    Soient $A,B\in\Z[X]$, et $B\neq0$. Montrer que si $B$ est unitaire, alors $Q,R\in\Z[X]$.
    \tcblower
    Même preuve que précedemment.
\end{exi}

\vspace*{-0.7cm}

\begin{ex}{}{}
    Faire la division euclidienne de $X^4+3X^3+X^2-X+1$ par $X^2+3X-4$.
    \tcblower
    Sans détailler :
    \begin{equation*}
        X^4+3X^3+X^2-X+1=(X^2+5)(X^2+3X-4)-16X+21
    \end{equation*}
\end{ex}

\begin{prop}{}{}
    Soit $P\in\K[X]$ et $a\in\K$.
    \begin{equation*}
        P(a) = 0 \iff (X-a) \mid P.
    \end{equation*}
    \tcblower
    \boxed{\ra} Par division euclidienne : $\exists (Q,R)\in\K[X] ~ P=(X-a)Q+R$, et $\deg(R)<1$ donc $R$ constant.\\
    On a alors $P(a)=R(a)=0$ donc $R=0$, donc $P=(X-a)Q$ donc $(X-a) \mid P$.\\ 
    \boxed{\la} $\exists Q \in \K[X], ~ P=(X-a)Q$, puis $P(a)=0\cdot Q(a)=0$.
\end{prop}

\vspace*{-0.7cm}

\begin{ex}{}{}
    Trouver tous les polynômes $P\in\R[X]$ tels que $(X+4)P(X)=XP(X+1)$.
    \tcblower
    \bf{Analyse.} Soit $P\in\R[X]$ tel que $(X+4)P(X)=XP(X+1)$.\\
    On a alors $P(0)=0$ et $P(-3)=0$, ainsi : $\exists Q_1 \in \R[X], ~ P=X(X+3)Q_1$.\\
    Alors $(X+4)(X+3)XQ_1(X)=X(X+1)(X+4)Q_1(X+1)$. On peut simplifier car $\K[X]$ est un anneau intègre.\\
    Alors $(X+3)Q_1(X)=(X+1)Q_1(X+1)$, puis $Q_1(-1)=Q_1(-2)=0$ : $\exists Q_2\in\K[X], ~ Q_1=(X+1)(X+2)Q_2$.\\
    Donc $(X+3)(X+2)(X+1)Q_2(X)=(X+3)(X+2)(X+1)Q_2(X+1)$ donc $Q_2(X)=Q_2(X+1)$.\\
    On en déduit que $Q_2$ est constant. Donc $P=\l X(X+1)(X+2)(X+3)$, avec $\l\in\K$.\\
    \bf{Synthèse.} On vérifie facilement que ces polynômes conviennent.
\end{ex}

\vspace*{-0.7cm}

\begin{thm}{Anneau principal $\star$.}{}
    Soit $I$ un idéal non nul de $\K[X]$. Alors $I$ est engendré par un unique polynôme unitaire.
    \begin{equation*}
        \exists P_0 \in \K[X] ~\nt{unitaire}, ~ I=(P_0).
    \end{equation*}
    Dans ce cas, $I$ est dit principal.
    \tcblower
    Soit $I$ un idéal non nul de $\K[X]$.\\
    \bf{Existence.} Notons $E=\{\deg(P), ~ P\in I, ~ P\neq 0\}$, sous-ensemble non-vide de $\N$. On note $n_0:=\min(E)$.\\
    Soit $P_0\in I$ tel que $\deg(P_0)=n_0$. On suppose $P_0$ unitaire SPDG, montrons que $I=(P_0)$.\\
    \boxed{\subset} Soit $P\in I$. Par division euclidienne : $\exists (Q,R)\in\K[X], ~ P=QP_0+R$ et $\deg(R)<\deg(P_0)$.\\
    On a $R=P-P_0Q\in I$, car c'est un groupe, donc $R=0$ par définition de $P_0$, puis $P=P_0Q\in (P_0)$.\\
    \boxed{\supset} On a $P_0\in I$ et $I$ est un idéal donc $I \supset (P_0)$.\n
    \bf{Unicité.} Supposons $I=(P)=(Q)$ avec $P$ et $Q$ unitaires.\\
    Alors $(P)\subset(Q)$ : $Q \mid P$, puis $(Q)\subset(P)$ : $P\mid Q$, donc $P$ et $Q$ sont associés. Ils sont unitaires, donc égaux.
\end{thm}

\vspace*{-0.4cm}

\bf{Remarque.} La preuve est en réalité la même que celle sur les sous-groupes de $\Z$.

\begin{defi}{PGCD.}{}
    Soient $(A,B)\in\K[X]^2$ non tous les deux nuls.\\
    Alors $I=A\K[X]+B\K[X]$ est un idéal non nul de $\K[X]$ engendré par un unique polynôme unitaire $D$.\\
    On appelle alors $D$ le PGCD de $A$ et $B$, noté $A\land B$.
    \tcblower
    On remarque que $I=(A)+(B)$, c'est donc un idéal de $\K[X]$ comme somme d'idéaux, non nul car $A,B\in I$.\\
    On conclut avec le théorème précédent.
\end{defi}

\vspace*{-0.7cm}

\begin{prop}{}{}
    Soient $A,B\in\K[X]$ non tous les deux nuls, on note $D=A\land B$.\\
    $\hra$ $D \mid A$ et $D \mid B$.\\
    $\hra$ Pour $C\in\K[X]$, si $C \mid A$ et $C\mid B$, alors $C \mid D$.\\
    Le PGCD est donc bien le plus grand diviseur commun de $A$ et $B$ au sens de la relation divise.
\end{prop}

\vspace*{-0.7cm}

\begin{prop}{Bézout. $\star$}{b1}
    Soient $A,B\in\K[X]$ non tous les deux nuls et $D=A\land B$.
    \begin{equation*}
        \exists (U,V) \in \K[X]^2 \quad D=AU+BV.
    \end{equation*}
\end{prop}

\vspace*{-0.7cm}

\begin{defi}{PPCM.}{}
    Soient $A,B\in\K[X]$ non nuls.\\
    Alors $I=A\K[X]\cap B\K[X]$ est un idéal non-nul de $\K[X]$ engendré par un unique polynôme unitaire $M$.\\
    On appelle alors $M$ le PPCM de $A$ et $B$, noté $A\lor B$.
    \tcblower
    On a $I=A\K[X]\cap B\K[X]$ un idéal de $\K[X]$ comme intersection d'idéaux de $\K[X]$, non nul car $AB\in I$.
\end{defi}

\vspace*{-0.7cm}

\begin{prop}{}{}
    Soient $A,B\in\K[X]$ non nuls, notons $M=A\lor B$.\\
    $\hra$ $A\mid M$ et $B\mid M$.\\
    $\hra$ Pour $N\in\K[X]$, si $A\mid N$ et $B\mid N$, alors $M \mid N$.\\
    Le PPCM est donc bien le plus petit multiple commun de $A$ et $B$ au sens de la relation divise.
\end{prop}

\vspace*{-0.7cm}

\subsection{Polynômes premiers entre eux.}

\begin{prop}{Caractérisation par Bézout.}{}
    Soit $(A,B)\in\K[X]^2$.
    \begin{equation*}
        A\land B = 1 \iff \exists (U,V) \in \K[X]^2, ~ AU+BV=1.
    \end{equation*}
    \tcblower
    \boxed{\ra} Égalité de Bézout (\ref{prop:b1}) avec $D=1$.\\
    \boxed{\la} Soit $D=A\land B$, alors $D\mid A$ et $D\mid B$, donc $D\mid AU+BV=1$ donc $D=1$.
\end{prop}

\vspace*{-0.8cm}

\begin{thm}{Caractérisation du PGCD. $\star$}{}
    Soient $A,B\in\K[X]$ non nuls et $D\in\K[X]$.
    \begin{equation*}
        D=A\land B \iff \exists (A',B')\in\K[X]^2, ~ (A = DA', ~ B=DB', ~ A'\land B'=1)
    \end{equation*}
    \tcblower
    \boxed{\ra} Alors $\exists (A',B')\in\K[X], ~ A=DA' \et B=DB'$ puis $\exists (U,V)\in\K[X]^2, ~ AU+BV=1$.\\
    Enfin, $D(A'U+B'V)=D$, on simplifie par $D$, il reste $A'U+B'V=1$, i.e. $A'\land B'=1$.\n
    \boxed{\la} On a $D$ diviseur commun de $A$ et $B$, puis $\exists (U,V)\in\K[X]^2, ~ A'U+B'V=1$, puis $AU+BV=D$.\\
    Alors $C\mid A$ et $C\mid B \ra C \mid D$ donc $D=A\land B$.
\end{thm}

\vspace*{-0.4cm}

\begin{prop}{}{}
    \begin{equation*}
        \left[A\land B=1 ~\et~ A\land C = 1\right] \ra A\land(BC)=1.
    \end{equation*}
    On le généralise avec un produit de $n$ polynômes.
    \tcblower
    \boxed{1.} Par Bézout : $AU_1+BV_1=1$ et $AU_2+CV_2=1$ donc par produit : $(AU_1+BV_1)(AU_2+CV_2)=1$.\\
    Ainsi : $A(AU_1U_2 + CU_1V_2 + BU_2V_1) + BC(V_1V_2) = 1$, donc par Bézout : $A\land (BC)=1$.\\
    \boxed{2.} Généralisation immédiate, par récurrence sur $n$.
\end{prop}

\vspace*{-0.8cm}

\begin{prop}{}{}
    \begin{equation*}
        \left[A\mid BC ~\et~ A\land B = 1\right]\ra A \mid C
    \end{equation*}
    \tcblower
    On a $\exists (U,V)\in\K[X]^2, ~ AU+BV=1$, d'où $ACU+(BC)V=C$, or $A\mid BC$ donc $A\mid C$. 
\end{prop}

\vspace*{-0.8cm}

\begin{prop}{}{}
    \begin{equation*}
        \left[A\mid C ~\et~ B\mid C ~\et~ A\land B=1\right] \ra AB\mid C.
    \end{equation*}
    \tcblower
    On a $\exists (U,V)\in\K[X]^2, ~ C=AU=BV$, alors $A\mid BV$ et $A\land B = 1$ d'où $A\mid V$.\\
    On a $\exists W \in \K[X], ~ V=AW$, puis $C=(AB)W$ donc $AB\mid C$.
\end{prop}

\vspace*{-0.8cm}

\begin{ex}{}{}
    Soit $P$ un polynôme admettant $a_1,...,a_n$ pour racines distinctes.\\
    Montrer que $(X-a_1)...(X-a_n)\mid P$.
    \tcblower
    Déjà, pour $a,b\in\K[X], ~ a\neq b \ra (X-a)\land(X-b)=1$. On a $(X-a_1)\mid P$,...$(X-a_n)\mid P$.\\
    Puis $(X-a_1)$,...,$(X-a_n)$ sont premiers deux-à-deux, donc $(X-a_1)...(X-a_n)\mid P$.
\end{ex}

\vspace*{-0.8cm}

\begin{corr}{}{}
    \begin{enumerate}
        \item Un polynôme de degré $n\in\N$ de $\K[X]$ admet au plus $n$ racines distinctes.
        \item Si un polynôme admet plus de racines distinctes que son degré, alors il est nul.
    \end{enumerate}
    \tcblower
    \boxed{1.} Soit $P$ de degré $n$, et $(a_1,...,a_r)$ ses racines distinctes : $(X-a_1)...(X-a_r)\mid P$, puis $r\leq n$ par degré.\\
    \boxed{2.} Par contraposée, soit $P$ non nul, de degré $n\in\N$, alors $P$ admet un nombre fini $r\leq n$ de racines.
\end{corr}

\vspace*{-0.8cm}

\begin{ex}{}{}
    Soit $P\in\R[X]$ et $a<b$ tels que $\forall x \in [a,b], ~ P(x)=0$, alors $P=0_{\K[X]}$.
    \tcblower
    Notons $n\in\N$ le degré de $P$, et soit $k > n$, on pose $\forall i \in \lb1,k\rb, ~ a_i = a + i\frac{b-a}{k}\in[a,b]$.\\
    Les $a_i$ sont tous distincts et dans $[a,b]$, donc $\forall i \in \lb1,k\rb, ~ P(a_i)=0$, donc $P$ admet $k>n$ racines distinctes.\\
    Par la proposition précédente, $P=0_{\K[X]}$.  
\end{ex}

\begin{ex}{}{}
    Soit $P=a_0+...+a_nX^n\in\K[X]$ polynôme formel. On note $\tilde{P}:x\mapsto a_0+...+a_nx^n$.\\
    De manière réciproque, si on connaît une fonction polynomiale, est-ce qu'on peut retrouver de manière unique le polynôme formel associé ? i.e. $\tilde{P}=\tilde{Q}\ra P=Q$ ?\\
    Si $\K$ est infini, la réponse est oui, car $\tilde{P}-\tilde{Q}$ admet une infinité de racines.\\
    Dans $\R[X]$ et $\C[X]$ on pourra confondre (via isomorphisme) polynômes formels et fonctions polynomiales.\\
    Par exemple dans $\Z/2\Z$, et $P=X^2-X$, on a $\tilde{P}=0$ et $P\neq0$.
\end{ex}

\vspace*{-0.8cm}

\subsection{Généralisation du PGCD, du PPCM.}

\begin{defi}{PGCD généralisé. $\star$}{}
    Soit $n\in\N^*\setminus\{1\}$, et $A_1,...,A_n\in\K[X]$ non tous nuls.
    \begin{equation*}
        I=(A_1)+...+(A_n)=\{A_1U_1+...+A_nU_n; ~ U_1,...,U_n\in\K[X]\}.
    \end{equation*}
    Alors $I$ est un idéal non nul de $\K[X]$ donc engendré par un unique polynôme unitaire $D$.
    \begin{center}
        $D$ est le PGCD de $A_1,...,A_n$, noté $A_1\land ... \land A_n$. 
    \end{center}
    En particulier, $\exists (U_1,...,U_n)\in\K[X], ~ D=A_1U_1+...+A_nU_n$.
\end{defi}

\vspace*{-0.8cm}

\begin{defi}{PPCM généralisé.}{}
    Soit $n\in\N^*\setminus\{1\}$, et $A_1,...,A_n\in\K[X]$ tous non nuls.
    \begin{equation*}
        I=(A_1)\cap...\cap(A_n) \nt{ est un idéal non-nul car } A_1...A_n\in I.
    \end{equation*}
    On a $I$ engendré par un unique polynôme unitaire $M$, et $M$ est le PPCM de $A_1,...,A_n$, noté $A_1\lor...\lor A_n$.
\end{defi}

\vspace*{-0.8cm}

\begin{defi}{Polynômes premiers entre-eux.}{}
    Soit $n\in\N^*\setminus\{1\}$, et $A_1,...,A_n\in\K[X]$ tous non nuls.\\
    Alors $A_1,...,A_n$ sont premiers entre-eux dans leur ensemble ssi $A_1\land...\land A_n=1$.
\end{defi}

\vspace*{-0.8cm}

\begin{prop}{Caractérisation par Bézout.}{}
    Soit $n\in\N^*\setminus\{1\}$, et $A_1,...,A_n$ tous non nuls. Alors :
    \begin{equation*}
        A_1\land ... \land A_n = 1 \iff \exists (U_1,...,U_n)\in\K[X], ~ A_1U_1+...+A_nU_n=1
    \end{equation*}
\end{prop}

\vspace*{-0.8cm}

\subsection{Polynômes irréductibles.}

\begin{defi}{Polynômes associés.}{}
    On dit que deux polynômes $P,Q\in\K[X]$ sont \bf{associés} ssi $P\mid Q$ et $Q\mid P$.
\end{defi}

\vspace*{-0.8cm}

\begin{defi}{}{}
    Soit $P\in\K[X]$ non-constant.\\
    On dit que $P$ est \bf{irréductible} ssi ses seuls diviseurs sont les polynômes constants et les polynômes associés à $P$.
\end{defi}

\vspace*{-0.8cm}

\begin{thm}{d'Alembert-Gauss.}{}
    Tout polynôme non constant de $\C[X]$ admet au moins une racine.
\end{thm}

\vspace*{-0.8cm}

\begin{corr}{$\star$}{}
    Les irréductibles de $\C[X]$ sont exactement les polynômes de degré $1$.
\end{corr}

\vspace*{-0.7cm}

\begin{corr}{}{}
    Tout polynôme de $\C[X]$ est scindé.
\end{corr}

\vspace*{-0.7cm}

\begin{thm}{Polynômes irréductibles de $\R$. $\star$}{}
    Les polynômes irréductibles de $\R[X]$ sont exactement :\\
    $\hra$ Les polynômes de degré 1.\\
    $\hra$ Les polynômes de degré 2 à discriminant négatif.
    \tcblower
    Il est clair que les polynômes cités sont irréductibles. Réciproquement, par contraposée, soit $P\in\R[X]$.\\
    $\hra$ Si $\deg(P)\geq3$, alors $P$ admet une racine $\a$ dans $\C$. Donc $\ov{\a}$ est aussi racine.\\
    Si $\a\in\R$, $(X-\a)\mid P$, sinon $(X-\a)(X-\ov{\a})=X^2-2\Re(\a)X+|\a|^2\mid P$.\\
    $\hra$ Si $\deg(P)=2$, on suppose qu'il admet une racine réelle, il admet un diviseur.
\end{thm}

\vspace*{-0.7cm}

\begin{prop}{}{}
    Soit $P\in\K[X]$ irréductible, et $A,B\in\K[X]$.
    \begin{equation*}
        P\mid AB \ra P\mid A ~\ou~ P\mid B
    \end{equation*}
    \tcblower
    Supposons $P\mid AB$ et $P\not| ~A$, alors $P\land A=1$ car $P$ irréductible, donc $P\mid AB$ et $P\land A = 1$, donc $P\mid B$. 
\end{prop}

\vspace*{-0.7cm}

\begin{thm}{Décomposition en polynômes irréductibles.}{}
    Soit $P\in\K[X]$ tel que $\deg(P)\geq1$. Alors $P$ s'écrit de manière unique sous la forme :
    \begin{equation*}
        P=\l P_1^{\a_1} \cdots P_r^{\a_r}, \quad r\in\N^*, ~ P_1,...,P_r ~ \nt{irréductibles}, ~ \a_1,...,\a_r \in \N^*.
    \end{equation*}
\end{thm}

\vspace*{-0.7cm}

\begin{prop}{}{}
    Soient $P_1,...,P_r\in\K[X]$ irréductibles, distincts et unitaires et :
    \begin{equation*}
        A = \l\prod_{i=1}^rP_i^{\a_i}; \quad B = \mu\prod_{i=1}^rP_i^{\a_i}.
    \end{equation*}
    Alors $A\land B = P_1^{\min(\a_1,\b_1)}\cdots P_r^{\min(\a_r,\b_r)}$ et $A\lor B = P_1^{\max(\a_1,\b_1)}\cdots P_r^{\max(\a_r,\b_r)}$.
\end{prop}

\vspace*{-0.7cm}

\begin{ex}{}{}
    Décomposer les polynômes suivantes en irréductibles dans $\R[X]$ et $\C[X]$ :
    \begin{enumerate}
        \item $X^3-1$.
        \item  $X^4-1$.
        \item  $X^n-1$.
        \item  $X^4+X^2+1$.
    \end{enumerate}
    \tcblower
    \boxed{1.} $X^3-1=(X-1)(X^2+X+1)=(X-1)(X-j)(X-j^2)$.\\
    \boxed{2.} $X^4-1=(X-1)(X+1)(X^2+1)=(X-1)(X+1)(X+i)(X-i)$.\\
    \boxed{3.} $X^n-1=\prod_{k=0}^{n-1}(X-e^{\frac{2ik\pi}{n}})=(X-1)\prod_{k=1}^{\frac{n-1}{2}}(X-e^{2ik\pi/n})(X-e^{-2ik\pi/n})$.\\
    \boxed{4.} $X^{4}+X^2+1=(X^2-j)(X^2-j^2)(X^2+j)(X^2+j^2)$.
\end{ex}    

\section{Un peu d'algèbre linéaire.}

\begin{defi}{Famille libre.}{}
    Soit $E$ un $\K$-ev et $(a_n)_{n\in\N}$ une famille infinie de vecteurs de $E$.
    \begin{equation*}
        (a_n)_{n\in\N} \nt{ est libre } \iff \nt{ toute sous-famille finie de $(a_n)_{n\in\N}$ est libre.} 
    \end{equation*}
\end{defi}

\vspace*{-0.7cm}

\begin{prop}{}{}
    Une famille de polynômes tous non nuls et de degrés distincts est libre.  
\end{prop}

\vspace*{-0.7cm}

\begin{ex}{Base de Taylor.}{}
    \begin{enumerate}
        \item Soit $a\in\K$, et $n\in\N^*$. Justifier que $\B=(1,(X-a),...,(X-a)^n)$ est une base de $\K_n[X]$.
        \item Donner les coordonnées de $P\in\K_n[X]$ dans cette base.
        \item Résoudre $P' \mid P$ dans $\K[X]$.
        \item Rappeler la caractérisation des racines multiples avec la formule de Taylor.
    \end{enumerate}
    \tcblower
    \boxed{1.} Déjà, c'est une famille de $\K_n[X]$ : $\forall k \in \lb0,n\rb, ~ (X-a)^k \in \K_n[X]$.\\
    De plus, cette famille est libre puisque ses polynômes sont de degrés tous distincts.\\
    Ainsi, $\B$ est une famille libre de $\K_n[X]$, de cardinal $n+1=\dim(\K_n[X])$ : c'est donc une base de $\K_n[X]$.\\
    \boxed{2.} Soit $P\in\K_n[X]$, on l'écrit dans la base $\B$ :
    \begin{equation*}
        \exists (\l_0,...,\l_n) \in \K^{n+1}, ~ P(X) = \sum_{k=0}^n\l_k(X-a)^k.
    \end{equation*} 
    En dérivant cette somme (finie) $m$ fois, pour $m \leq n$, et en évaluant en $a$, on obtient :
    \begin{equation*}
        P^{(m)}(a) = m!\l_m + \sum_{k=1}^{n-m}\l_{k+m}(a-a)^k = m!\l_m \ra \l_m = \frac{P^{(m)}(a)}{m!}.
    \end{equation*}
    \boxed{3.} Si $P'\mid P$, alors $\exists Q \in \K[X], ~ P = P'Q$.\\
    $\hra$ Par passage au degré, $\deg(Q)=1$, donc $\exists \l,a \in \K^*, ~ Q = \l(X-a)$.\\
    Ainsi, $P=\l(X-a)P'$ puis $\l = \frac{1}{n}$ par égalité des coefficients dominants, puisque $\cd(P)=a_n$ et $\cd(P')=na_n$.\\
    Dans la base de Taylor :
    \begin{equation*}
        P = \sum_{k=0}^n\l_k(X-a)^k; \et \l(X-a)P' = \frac{1}{n}\sum_{k=0}^nk\l_k(X-a)^{k} \ra \forall k \in \lb0,n\rb, ~ \l_k = \frac{k}{n}\l_k.
    \end{equation*}
    Donc si $k<n, ~ \l_k = 0$, ce qui laisse $P=\l_n(X-a)^n$.
\end{ex}

\vspace*{-0.7cm}

\begin{ex}{}{}
    Soit $P=(X-1)^{31}(X-2)^2$. Calculer $P^{(n)}(1)$ pour $n\in\N$.
    \tcblower
    $\hra$ $\deg(P)=33$, donc $\forall n > 33, ~ P^{(n)}=0_{\R[X]}$ donc $P^{(n)}(1)=0$.\\
    $\hra$ 1 est racine de multiplicité $31$ donc $\forall n \in \lb0,30\rb, ~ P^{(n)}(1)=0$.\\
    $\hra$ Taylor : $(X-2)^2=(X-1)^2-2(X-1)+1$ donc $P=(X-1)^{33}-2(X-1)^{32}+(X-1)^{31}$.\\
    Par identification : $P^{(31)}(1)=31!$, $P^{(32)}(1)=-2\times32!$, $P^{(33)}(1)=33!$.
\end{ex}

\vspace*{-0.7cm}

\begin{prop}{Polynômes échelonnés. $\star$}{}
    \begin{enumerate}
        \item Soit $n\in\N^*$ et $(P_k)_{k\leq n}$ une famille \bf{finie} de polynômes échelonnés de $\K_n[X]$, alors $(P_k)$ est base de $\K_n[X]$.
        \item Soit $(P_n)_{n\in\N}$ une famille \bf{infinie} échelonnée de $\K[X]$. Alors $(P_n)_{n\in\N}$ est une base de $\K[X]$.
    \end{enumerate}
    \tcblower
    \boxed{1.} $(P_k)$ est une famille de polynômes non nuls de degrés distincts donc libre.\\
    Notons $\B=(P_0,...,P_n)$, c'est une famille libre de $\K_n[X]$, et $|\B|=n+1=\dim(\K_n[X])$ donc $\B$ est une base.\\
    \boxed{2.} Par les mêmes arguments, $(P_n)$ est une famille libre de $\K[X]$ car toutes ses familles finies sont libres.\\
    Soit $P\in\K[X]$. Si $P=0$, trivial. Sinon, on note $n=\deg(P)\in\N$. Alors $P\in\K_n[X]$.\\
    Cependant, $(P_0,...,P_n)$ est base de $\K_n[X]$, donc $P\in\Vect(P_1,...,P_n)$, donc $(P_n)_{n\in\N}$ est base de $\K[X]$.
\end{prop}

\vspace*{-0.8cm}

\begin{exercice}{Polynômes de Lagrange. $\star$}{}
    Soient $a_0<...<a_n\in\R$. On définit $L_i\in\K_n[X]$ par $L_i(a_j)=\d_{i,j}$ pour $i,j\in\lb0,n\rb$.
    \begin{enumerate}
        \item Vérifier que $L_i$ est défini de manière unique et donner une expression explicite de $L_i$.
        \item Montrer que $(L_0,...,L_n)$ est une base de $\K_n[X]$.
        \item Donner les coordonnées de $P\in\K_n[X]$ dans cette base.
        \item En déduire que $\forall (b_0,...,b_n)\in\K^{n+1}, ~ \exists!P\in\K_n[X], ~ \forall i \in \lb0,n\rb$ tel que $P(a_i)=b_i$.
        \item Retrouver le résultat de $(4)$ avec $\phi:P\mapsto(P(a_0),...,P(a_n))$.
        \item Soit $p\in\N$. Calculer $\sum_{i=0}^nL_i$, $\sum_{i=0}^na_i^pL_i$.
    \end{enumerate}
    \tcblower
    \small
    \boxed{1.} \bf{Analyse.} Soit $i\in\lb0,n\rb$, alors $\forall j \neq i, ~ L_i(a_j) = 0$, donc :
    \begin{equation*}
        \prod_{\substack{j=0 \\ j\neq i}}^n (X-a_j) ~ \mid ~ L_i ~~\nt{alors}~~ \exists Q \in \K[X], ~ L_i = Q(X)\prod_{\substack{j=0\\j\neq i}}^n(X-a_j).
    \end{equation*}
    Par passage au degré, $\deg(Q)=0$ : $\exists \l \in \K, ~ Q = \l$, et $L_i(a_i) = 1$, donc :
    \begin{equation*}
        L_i(a_i) = \l \prod_{j\neq i}(a_i - a_j) \ra \l = \frac{1}{\prod_{j\neq i}(a_i - a_j)}.
    \end{equation*}
    Cette division a un sens car les $(a_j)$ sont distincts deux-à-deux. On a donc obtenu l'unicité de $L_0$.\\
    \bf{Synthèse.} On vérifie facilement que $\deg(L_i)=n$ et que $\forall j \in \lb0,n\rb, ~ L_i(a_j)=\d_{i,j}$.\\
    \boxed{2.} Montrons que la famille $(L_0,...,L_n)$ est libre dans $\K_n[X]$.\\
    $\hra$ Déjà, c'est bien une famille de $\K_n[X]$ puisque $\forall i \in \lb0,n\rb, ~ \deg(L_i)=n$.\\
    $\hra$ Soient $(\l_0,...,\l_n)\in\K^{n+1}$ tels que $\l_0L_0+...+\l_nL_n=0_{\K[X]}$.\\
    Alors en évaluant en $a_i$ pour tout $i\in\lb0,n\rb$, on obtient $\l_i = 0$ : c'est bien une famille libre.\\
    C'est donc une famille libre de $\K_n[X]$ de cardinal $n+1=\dim(\K_n[X])$ : c'est une base de $\K_n[X]$.\\
    \boxed{3.} Soit $P\in\K_n[X]$. Puisque $(L_0,...,L_n)$ est une base de $\K_n[X]$, $\exists (\l_0,...,\l_n)\in\K^{n+1}, ~ P=\l_0L_0+...+\l_nL_n$.
    \begin{equation*}
        \forall i \in \lb0,n\rb, ~ P(a_i) = \l_i \ra P = \sum_{i=0}^n P(a_i)L_i.
    \end{equation*}
    \boxed{4.} Soit $(b_0,...,b_n)\in\K^{n+1}$. Le polynôme suivant convient : (unicité facile)
    \begin{equation*}
        P = \sum_{i=0}^n b_iL_i ~ \nt{alors} ~ \forall i \in \lb 1, n\rb, ~ P(a_i) = b_iL_i(a_i) = b_i.
    \end{equation*}
    \boxed{5.} $\phi$ est une application linéaire de $\K_n[X]$ dans $\K^{n+1}$ injective ($\Ker(\phi)=\{0\}$) puis bijective par égalité des dimensions. À chaque $n$-uplet $(b_0,...,b_n)\in\K^{n+1}$ on peut donc associer un unique polynôme $P\in\K_n[X]$ qui convient.\\
    \boxed{6.} Soit $P\in\K_n[X]$, alors :
    \begin{equation*}
        P = \sum_{i=0}^nP(a_i)L_i \ra \sum_{i=0}^nL_i = 1 \et \sum_{i=0}^na_i^pL_i = X^p.
    \end{equation*}
\end{exercice}

\vspace*{-0.7cm}

\section{Fractions rationnelles.}

\begin{defi}{Fraction rationnelle.}{}
    Une fonction \bf{fraction rationnelle} est de la forme $x\mapsto~\frac{P(x)}{Q(x)}$ avec $P,Q\in\K[X]$.
\end{defi}

\vspace*{-0.7cm}

\begin{defi}{}{}
    On notera $\K(X)$ l'ensemble des fractions rationnelles, c'est plus petit corps qui contient $\K[X]$.
\end{defi}

\vspace*{-0.7cm}

\begin{prop}{}{}
    Soient $P,Q\in\K[X]$, avec $P\land Q = 1$. On a $\deg\left( \frac{P}{Q} \right)=\deg(P)-\deg(Q)\in\Z$.
\end{prop}

\vspace*{-0.7cm}

\begin{thm}{}{50}
    Soient $P,Q\in\K[X]$ non nuls. Alors $\exists!E\in\K[X], ~ \exists!G\in\K(X), ~ \frac{P}{Q}=E+G$ avec $\deg(G)<0$.\\
    Alors on appelle $E$ la partie entière, et $G$ la partie fractionnaire de $P/Q$.
    \tcblower
    Par division euclidienne : $\exists!(E,R)\in\K[X]\times\K(X), ~ P=QE+R$ avec $\deg(R)<\deg(Q)$.\\
    Puis on a $\frac{P}{Q}=E+\frac{R}{Q}$, on pose $G:=\frac{R}{Q}$, on a alors $\deg(G)=\deg(R)-\deg(Q)<0$.
\end{thm}

\begin{defi}{Racines, pôles.}{}
    Soient $P,Q\in\K[X]$ non nuls et tels que $P\land Q = 1$. Soit $F=P/Q$.
    \begin{itemize}
        \item Les \bf{racines} de $F$ (comptées avec leurs multiplicité) sont les racines de $P$.
        \item Les \bf{pôles} de $F$ (comptés avec leurs multiplicité) sont les racines de $Q$.
    \end{itemize}
\end{defi}

\begin{thm}{Calcul des coefficients associés aux pôles simples.}{}
    Soient $P,Q\in\K[X]$, et $a$ une racine simple de $Q$ : $\exists R \in \K[X], ~ Q=(X-a)R$ et $R(a)\neq0$.\\
    Si $P\land Q = 1$, alors :
    \begin{equation*}
        \exists!\l \in \K, ~ \frac{P}{Q} = \frac{\l}{X-a} + F
    \end{equation*}
    où $F\in\K(X)$, et $a$ n'est pas pôle de $F$.\\
    $\hra$ En multipliant par $X-a$ puis en évaluant en $a$, on trouve $\l=\frac{P(a)}{R(a)}$ (formule du cache).\\
    $\hra$ En remarquant que $Q'(a)=R(a)$, on a aussi $\l=\frac{P(a)}{Q'(a)}$ (formule de la dérivée).\\
    \bf{Remarque.} Si $P\land Q\neq 1$, revenir au théorème \ref{thm:50}.
\end{thm}

\begin{ex}{}{}
    Pour $\ds F=\frac{X+4}{(X-2)(X-3)}$, $\exists!(\l,\mu)\in\K^2, ~ F=\frac{\l}{X-2}+\frac{\mu}{X-3}$, puis $(\l,\mu)=(-6,6)$.
\end{ex}

\begin{ex}{}{}
    Pour $F=\frac{X^3+4}{(X-1)(X-2)}$, on calcule la partie entière : $X^3+4=(X+3)(X^2-3X+2)+7X-2$.\\
    On a alors $F=X+3+\frac{7X-2}{(X-1)(X-2)}=X+3+\frac{\l}{X-1}+\frac{\mu}{X-2}$, avec $(\l,\mu)=(-5,12)$.
\end{ex}

\vspace*{-0.8cm}

\begin{ex}{}{}
    Décomposer en éléments simples $F=\frac{1}{X^3-1}$ dans $\C[X]$ puis dans $\R[X]$.
    \tcblower
    $\bullet$ Dans $\C[X]$, cette fraction a trois pôles simples : $1,j,j^2$ : $\exists(\l,\mu,\g)\in\C^3, ~ F=\frac{\l}{X-1}+\frac{\mu}{X-j}+\frac{\g}{X-j^2}$.\\
    $\hra$ On multiplie par $(X-1)$ et on évalue en 1 : $\l = \frac{1}{(1-j)(1-j^2)}=\frac{1}{2-(j+j^2)}=\frac{1}{3}$.\\
    $\hra$ On multiplie par $(X-j)$ et on évalue en $j$ : $\mu = \frac{1}{(j-1)(j-j^2)}$.\\
    $\hra$ On multiplie par $(X-j^2)$ et on évalue en $j^2$ : $\g=\frac{1}{(j^2-1)(j^2-j)}$.\\
    $\bullet$ Dans $\R[X]$, on a $F=\frac{1}{(X-1)(X^2+X+1)}=\frac{\l}{X-1} + \frac{\mu X + \g}{X^2+X+1}$.\\
    On retrouve $\l=\frac{1}{3}$ de la même manière, puis en évaluant en $0$ : $-1 = -\frac{1}{3} + \g$ donc $\g = -\frac{2}{3}$.\\
    En multipliant par $(X^2+X+1)$ et en évaluant en -1, on trouve $\mu = -\frac{1}{3}$.
\end{ex}

\vspace*{-0.8cm}

\begin{exercice}{}{}
    Décomposer en éléments simples dans $\C[X]$ la faction rationnelle suivante :
    \begin{equation*}
        \frac{1}{X^n-1}.
    \end{equation*}
    \tcblower
    Soit $\w=e^{\frac{2i\pi}{n}}$, alors $\forall k \in \N, ~ \w^k$ est pôle simple de la fraction rationnelle :
    \begin{equation*}
        \exists (l_0,...,l_n) \in \K^{n+1}, ~ \frac{1}{X^n - 1} = \sum_{k=0}^n \frac{\l_k}{X-\w^k}.
    \end{equation*}
    Par la formule de la dérivée : $\l_k = \frac{1}{n\w^{k(n-1)}}=\frac{1}{nw^{kn}\w^{-k}}=\frac{\w^k}{n}$ car $w^{kn}=1$. On conlut :
    \begin{equation*}
        \frac{1}{X^n-1} = \frac{1}{n}\sum_{k=0}^n\frac{\w^k}{X-w^k}.
    \end{equation*}
\end{exercice}

\vspace*{-0.8cm}

\begin{ex}{}{}
    Soit $P\in\K[X]$ scindé, avec $n=\deg(P)\geq1$. Décomposer $\frac{P'}{P}$ en éléments simples.
    \tcblower
    Soit $P=\l(X-a_1)^{\a_1}...(X-a_n)^{\a_n}$, avec $a_1,...,a_n$ distincts, $(\a_1,...,\a_n)\in\N^*$ et $\l\in\K^*$.\\
    $\hra$ $\deg(\frac{P'}{P})=-1$, donc la partie entière de $\frac{P'}{P}$ est nulle.\\
    Soit $i\in\lb1,n\rb$ alors $\exists Q \in \K[X], ~ P(X) = (X - a_i)^{\a_i}Q(X)$ et $Q(a_i)\neq0$, donc :
    \begin{equation*}
        \frac{P'(X)}{P(X)} = \frac{\a_i(X-a_i)^{\a_i-1}Q(X) + (X-a_i)^{\a_i}Q'(X)}{(X-a_i)^{\a_i}Q(X)} = \frac{\a_i}{X-a_i} + \frac{Q'(X)}{Q(X)}.
    \end{equation*}
    où $a_i$ n'est pas pôle de $\frac{Q'(X)}{Q(X)}$, donc :
    \begin{equation*}
        \frac{P'(X)}{P(X)} = \sum_{i=1}^n \frac{\a_i}{X-a_i}.
    \end{equation*}
\end{ex}

\begin{ex}{}{}
    Décomposer la fraction rationnelle suivante en éléments simples :
    \begin{equation*}
        F = \frac{1}{(X-2)^3(X+3)}.
    \end{equation*}
    \tcblower
    $\hra$ $\deg(F)<0$ donc la partie entière de $F$ est nulle. Donc :
    \begin{equation*}
        F = \frac{\a}{X+3} + \frac{\b_1}{X-2} + \frac{\b_2}{(X-2)^2} + \frac{\b_3}{(X-2)^3}.
    \end{equation*}
    Puis par la formule du cache, on trouve $\a=-\frac{1}{5^3}$ et $\b_3 = \frac{1}{5}$.\\
    $\hra$ Posons $F_1 = F - \frac{\b_3}{(X-2)^3}=-\frac{1}{5(X-2)^2(X+3)}$; mais aussi $F_1=\frac{\a}{X+3}+\frac{\b_1}{X-2}+\frac{\b_2}{(X-2)^2}$.\\
    Par la formule du cache, on trouve $\b_2=-\frac{1}{5^2}$.\\
    $\hra$ Posons $F_2 = F_1 - \frac{\b_2}{(X-2)^2}=\frac{1}{5^2(X-2)(X+3)}$; mais aussi $F_2=\frac{\a}{X+3}+\frac{\b_1}{X-2}$.\\
    Par la formule du cache, on trouve $\b_1=\frac{1}{5^3}$.
    \begin{equation*}
        F = -\frac{1}{125(X+3)} + \frac{1}{125(X-2)} - \frac{1}{25(X-2)^2} + \frac{1}{5(X-2)^3}.
    \end{equation*}
\end{ex}

\begin{defi}{Nombres algébriques, transcendants. (HP)}{}
    Soit $\a\in\R$ et $I=\{P\in\Q[X], ~ P(\a)=0\}$ un idéal de $(\Q[X],+,\times)$.\\
    Si $I=(0)$, on dit que $\a$ est transcendant, sinon $I\neq(0)$ et $\exists!\pi\in\Q[X]$ unitaire tel que $I=(\pi)$.\\
    On dit alors que $\a$ est algébrique, et $\pi$ est appelé polynôme minimal de $\a$.\\
    Alors $\pi$ est aussi le polynôme de plus bas degré (non nul) dans $\Q[X]$ qui admet $\a$ comme racine.\\
    Par définition, le degré de $\a$ est le degré de $\pi$.\\
    \bf{Remarque.} Le polynôme minimal, s'il existe, est forcément irréductible dans $\Q[X]$.
\end{defi}

\begin{ex}{}{}
    Si $\a\in\Q$, alors $P=X-\a\in\Q[X]$, et $P(\a)=0$, $I=(P)$, $\pi=X-\a$ et $\a$ est de degré 1.\\
    Réciproquement, si $\a$ est algébrique de degré 1, alors $\a\in\Q$.
\end{ex}

\begin{ex}{}{}
    Pour $\a=\sqrt{2}$, on a $P=X^2-2\in\Q[X]$ annule $\a$, donc $\sqrt{2}$ est algébrique.
\end{ex}

\fin

\pagebreak

\section{Khôlles, exercices supplémentaires...}

\setexercicecounter{60}

\begin{exercice}{CCINP $\star$}{}
    Soit la matrice $A=\begin{pmatrix}1&2\\2&4\end{pmatrix}$ et $f$ l'endormorphisme de $M_2(\R)$ défini par $f(M)=AM$.
    \begin{enumerate}
        \item Déterminer une base de $\Ker f$.
        \item $f$ est-il surjectif ?
        \item Déterminer une base de $\Im f$.
        \item A-t-on $M_2(\R)=\Ker f\oplus \Im f$ ?
    \end{enumerate}
    \tcblower
    \boxed{1.} Soit $M=\begin{pmatrix}a&b\\c&d\end{pmatrix}$. On a $f(M)=\begin{pmatrix}a+2c&b+2d\\2a + 4c&2b+4d\end{pmatrix}$ d'où $M\in\Ker f \iff M = \begin{pmatrix}-2c & -2d \\ c & d\end{pmatrix}$.\\
    On pose $M_1 =\begin{pmatrix}-2 & 0 \\ 1 & 0\end{pmatrix}, M_2 = \begin{pmatrix}0 & -2 \\ 0 & 1\end{pmatrix}$, alors $\Ker f = \Vect\left(M_1, M_2\right)$ et $(M_1,M_2)$ est libre, c'est une base.\\
    \boxed{2.} On a $\rg(A)=1\neq2$, donc $A$ n'est pas inversible, donc $I_2$ n'a pas d'antécédent par $f$ : pas surjective.\\
    \boxed{3.} Théorème du rang : $\rg f = 4 - \dim \Ker f = 2$. La base $\left( \begin{pmatrix}1&0\\2&0\end{pmatrix}, \begin{pmatrix}0&2\\0&4\end{pmatrix} \right)$ convient.\\
    \boxed{4.} On a $\dim M_2(\R)=\dim \Ker f + \dim \Im f$. On montre facilement que $\Ker f \cap \Im f = \{0\}$.
\end{exercice}

\setexercicecounter{64}

\begin{exercice}{CCINP $\star$}{}
    Soit $f$ un endormorphisme d'un espace vectoriel $E$ de dimension finie $n$.
    \begin{enumerate}
        \item Démontrer que $E=\Im f \oplus \Ker f \ra \Im f = \Im f^2$.
        \item \begin{enumerate}
            \item Démontrer que $\Im f = \Im f^2 \iff \Ker f = \Ker f^2$.
            \item Démontrer que $\Im f = \Im f^2 \ra E = \Im f \oplus \Ker f$.
        \end{enumerate}
    \end{enumerate}
    \tcblower
    \boxed{1.} Supposons $E=\Im f \oplus \Ker f$. L'inclusion $\Im f \supset \Im f^2$ est triviale.\\
    \boxed{\subset} Soit $z\in \Im f$, $\exists (y,a) \in \Im f\times \Ker f, ~ z=f(y+a)=f(y)$ puis $\exists x \in E, ~ y=f(x)$ donc $z=f^2(x)$.\\
    \boxed{2.a)} Théorème du rang : $\Im f = \Im f^2 \iff \rg f = \rg f^2 \iff \dim \Ker f = \dim \Ker f^2$.\\
    Or $\Ker f \subset \Ker f^2$, donc par égalité des dimensions, $\Ker f = \Ker f^2$.\\
    \boxed{2.b)} Supposons $\Im f = \Im f^2$. On a $\Ker f = \Ker f^2$ d'après la question précédente.\\
    Soit $x\in \Ker f \cap \Im f$, on a $\exists a \in E, ~ f(a)=x, ~ f(x)=0$. Or $\Ker f = \Ker f^2$, d'où $x=f(a)=0$.\\
    Par théorème du rang, on a $E=\Im f + \Ker f$, d'où la somme directe.
\end{exercice}

\pagebreak

\setexercicecounter{85}

\begin{exercice}{CCINP $\star$}{}
    \begin{enumerate}
        \item Soient $n\in\N^*$, $P\in\R_n[X]$ et $a\in\R$.
        \begin{enumerate}
            \item Décomposer $P(X)$ dans la base $(1,X-a,(X-a)^2,...,(X-a)^n)$ avec la formule de Taylor.
            \item Soit $r\in\N^*$. En déduire que $a$ est racine de $P$ de multiplicité $r$ ssi $P^{(r)}(a)\neq0$\\ et $\forall k \in\lb0,r-1\rb, ~ P^{(k)}(a)=0$. 
        \end{enumerate}
        \item Déterminer deux réels $a$ et $b$ pour que $1$ soit racine double de $P=X^5+aX^2+bX$ et factoriser alors ce polynôme dans $\R[X]$.
    \end{enumerate}
    \tcblower
    \boxed{1.a)} $P(X)=\sum_{k=0}^n\frac{P^{(k)}(a)}{k!}(X-a)^k$.\\
    \boxed{1.b)} $a$ racine de multiplicité $r \iff (X-a)^r \mid P \iff \exists Q \in \R_{n-r}[X], ~ P=(X-a)^rQ$ et $Q(a)\neq0$.\\
    $\iff P=(X-a)^r\sum_{k=0}^{n-r}q_k(X-a)^k=\sum_{k=0}^{n-r}q_k(X-a)^{r+k}=\sum_{k=r}^{n}q_{k-r}(X-a)^k$.\\
    On voit bien que $\forall k \in \lb0,r-1\rb,~ P^{(k)}(a)=0$, et $P^{(r)}(a)\neq0$.\\
    \boxed{2.} $P(1)=P'(1)=0$ et $P''(1)\neq0$ donc $1+a+b=0$, $5+2a+b=0$ et $20+2a\neq0$.\\
    Alors $a+b=-1$ et $2a+b=-5$, puis $a=-1-b$ et $b=3$, puis $(a,b)=(-4,3)$.\\
    On a $X^5+aX^2+bX=X(X-1)^2(X^2+2X+3)$.
\end{exercice}

\setexercicecounter{87}

\begin{exercice}{CCINP $\star$}{}
    Soient $a_0,...,a_n\in\R$ deux-à-deux distincts.
    \begin{enumerate}
        \item Montrer que si $b_0,...,b_n\in\R$, alors existe $P\in\R[X]$ tel que $\deg(P)\leq n$ et $\forall i\in\lb0,n\rb, ~ P(a_i)=b_i$.
        \item Soit $i\in\lb0,n\rb$. Expliciter ce polynôme, que l'on notera $L_i$ lorsque :
        \begin{equation*}
            \forall j \in \lb0,n\rb, ~ b_j = \begin{cases}0 &\nt{si} \quad i\neq j \\ 1 &\nt{si} \quad i = j\end{cases}
        \end{equation*}
        \item Prouver que $\forall p \in \lb0,n\rb, ~ \sum_{k=0}^na_k^pL_k=X^p$.
    \end{enumerate}
    \tcblower
    \boxed{1.} Soit $\phi:P\mapsto (P(a_0),...,P(a_n))$, montrons que c'est un isomorphisme de $\R_n[X]$ dans $\R^{n+1}$.\\
    $\hra$ C'est bien une forme linéaire par linéarité de l'évaluation de polynômes.\\
    $\hra$ Soit $P\in\Ker\phi$, $P$ admet $n+1$ racines tout en étant de degré $n$, donc $P$ est nul, donc $\phi$ est injective.\\
    Par égalité des dimensions de $\R_n[X]$ et de $\R^{n+1}$, $\phi$ est donc aussi bijective.\\
    On a bien montré que $\forall (b_0,...,b_n)\in\R^{n+1}, ~ \exists!P\in\R_n[X], ~ \forall i \in \lb0,n\rb, ~ P(a_i)=b_i$.\\
    \boxed{2.} On reconnaît un polynôme d'interpolation de Lagrange, qui s'exprime :
    \begin{equation*}
        L_i = \prod_{\substack{j = 0\\j\neq i}}^n \frac{X-a_j}{a_i-a_j}.
    \end{equation*}
    C'est bien un polynôme de degré $n$, et on vérifie facilement les conditions sur les $b_j$.\\
    \boxed{3.} Soit $p\in\lb0,n\rb$ et $i\in\lb0,n\rb$, on évalue les deux polynômes en $a_i$ :
    \begin{equation*}
        \sum_{k=0}^na_k^pL_k(a_i) = \sum_{k=0}^na_k^p\d_{i,k} = a_i^p; \et X^p(a_i) = a_i^p.
    \end{equation*}
    Les deux polynômes coïncident donc sur les tous les $a_i$, donc par l'unicité vue en 1, ils sont égaux.
\end{exercice}

\end{document}